\documentclass[12pt,a4paper]{article}
\usepackage[utf8]{inputenc}
\usepackage[spanish,es-tabla]{babel}
\usepackage{amsmath,amsfonts,amssymb}
\usepackage{booktabs}
\usepackage{tabularx}
\usepackage{geometry}
\usepackage{graphicx}
\usepackage{float}
\usepackage{hyperref}
\usepackage{cite}

\geometry{top=2.5cm, bottom=2.5cm, left=3cm, right=3cm}

\hypersetup{
    colorlinks=true,
    linkcolor=black,
    citecolor=blue,
    urlcolor=blue
}

\begin{document}

% --- PORTADA ---
\begin{titlepage}
    \centering
    \LARGE \textbf{Universidad Tecnológica Nacional}\\
    \Large Facultad Regional Rosario\\
    \large Departamento de Ingeniería en Sistemas de Información\\[2.5cm]
    
    \Huge \textbf{Trabajo Práctico Integrador}\\[0.5cm]
    \LARGE Simulación del Flujo de Votantes y Tiempos de Espera en la Escuela de Educación Secundaria Técnica N°1 Fray Luis Beltrán de Arrecifes\\[2.5cm]
    
    \large
    \textbf{Cátedra:} Simulación\\
    \textbf{Docentes:} Juan Torres, Guillermo Leale\\[1.5cm]
    
    \textbf{Integrantes:}\\
    Manuel Casanovas\\
    Valentino Maseda\\
    Valentín Peppino\\
    Gianlucas Zabaleta\\[2cm]
    
    \vfill
    \large Rosario, 2026
\end{titlepage}

\tableofcontents
\newpage

% --- 1. INTRODUCCIÓN Y PLANIFICACIÓN ---
\section{Introducción y Planificación del Estudio}
El presente estudio de simulación aborda el análisis del flujo de electores y la congestión en el centro de votación emplazado en la Escuela de Educación Secundaria Técnica N°1 Fray Luis Beltrán de Arrecifes. En comicios de alta concurrencia, las mesas de votación suelen presentar demoras excesivas originadas en la etapa de selección de boletas dentro del cuarto oscuro, generando filas que interfieren con la circulación general de los pasillos escolares.

El objetivo general del trabajo consiste en evaluar el impacto operativo de implementar dos biombos de votación en paralelo por cada mesa electoral frente al esquema tradicional de biombo único. Como objetivos específicos se definen:
\begin{itemize}
    \item Modelar el comportamiento dinámico de los electores en el layout arquitectónico real del establecimiento.
    \item Medir las variables de desempeño del sistema: tiempo medio de permanencia total ($W$), demora promedio en fila de espera ($d(n)$) y longitud máxima de colas ($Q(t)$).
    \item Comprobar mediante contraste de hipótesis si la reducción temporal lograda con el esquema alternativo resulta estadísticamente significativa.
\end{itemize}

% --- 2. RECOLECCIÓN DE DATOS Y MODELO CONCEPTUAL ---
\section{Recolección de Datos y Definición del Modelo}
Para la caracterización del centro de votación, se relevaron los registros correspondientes a las 9 mesas electorales asignadas al establecimiento (mesas 38 a 46). El padrón electoral totalizó 3.146 electores habilitados, con una concurrencia efectiva de 2.199 votantes a lo largo de una jornada de 10 horas (08:00 a 18:00 hs), representando un 69,90\% de participación promedio (Tabla~\ref{tab:participacion_mesas}).

\begin{table}[H]
    \centering
    \small
    \setlength{\tabcolsep}{6pt}
    \caption{Padrón y nivel de participación por mesa electoral.}
    \label{tab:participacion_mesas}
    \begin{tabular}{ccccc}
        \toprule
        \textbf{Mesa} & \textbf{Electores en Padrón} & \textbf{Total Votantes} & \textbf{Ausentes} & \textbf{\% Participación} \\
        \midrule
        38 & 348 & 243 & 105 & 69,83\% \\
        39 & 350 & 235 & 115 & 67,14\% \\
        40 & 348 & 225 & 123 & 64,66\% \\
        41 & 350 & 255 & 95  & 72,86\% \\
        42 & 350 & 245 & 105 & 70,00\% \\
        43 & 351 & 239 & 112 & 68,09\% \\
        44 & 350 & 268 & 82  & 76,57\% \\
        45 & 349 & 258 & 91  & 73,93\% \\
        46 & 350 & 231 & 119 & 66,00\% \\
        \midrule
        \textbf{Total} & \textbf{3.146} & \textbf{2.199} & \textbf{947} & \textbf{69,90\%} \\
        \bottomrule
    \end{tabular}
\end{table}

\subsection{Parámetros de Servicio y Supuestos Estocásticos}
Los tiempos operativos asociados a cada etapa del circuito electoral fueron ajustados mediante distribuciones triangulares continuas expresadas en segundos, caracterizadas por sus valores mínimo, moda y máximo:
\begin{itemize}
    \item \textbf{Acreditación y verificación de identidad:} $\text{Triangular}(20, 35, 60)\text{ s}$. Proceso atendido por las autoridades de mesa (capacidad unitaria).
    \item \textbf{Votación en cuarto oscuro / biombo:} $\text{Triangular}(30, 60, 120)\text{ s}$. Representa el cuello de botella físico del sistema. Su capacidad varía según el escenario ($c = 1$ o $c = 2$).
    \item \textbf{Egreso, firma de padrón y entrega de constancia:} $\text{Triangular}(15, 25, 40)\text{ s}$. Cierre formal del sufragio por parte de las autoridades de mesa.
\end{itemize}

La tasa media de arribo global al establecimiento se dimensionó en $219,9\text{ votantes/hora}$, con un límite de arribos fijado en 2.199 entidades. La derivación hacia cada mesa se modeló mediante probabilidades de ruteo proporcionales a la concurrencia real observada en los comicios.

% --- 3. VALIDACIÓN DEL MODELO CONCEPTUAL ---
\section{Validación del Modelo Conceptual}
El circuito conceptual se estructuró a partir de una secuencia estricta de eventos discretos desacoplados en el espacio físico:
\begin{equation*}
    \text{Ingreso General} \xrightarrow{} \text{Desplazamiento a Mesa} \xrightarrow{} \text{Cola Mesa (FIFO)} \xrightarrow{} \text{Acreditación}
\end{equation*}
\begin{equation*}
    \xrightarrow{} \text{Desplazamiento a Biombo} \xrightarrow{} \text{Espera / Voto} \xrightarrow{} \text{Retorno a Mesa} \xrightarrow{} \text{Firma y Salida}
\end{equation*}

Se verificó conceptualmente la pertinencia del desacople entre la acreditación y el cuarto oscuro: al incorporar dos biombos, las autoridades pueden acreditar a un segundo elector mientras el primero selecciona sus boletas, reduciendo el tiempo ocioso y mitigando la cola exterior en el pasillo escolar.

% --- 4. CONSTRUCCIÓN Y VERIFICACIÓN DEL PROGRAMA ---
\section{Construcción y Verificación del Programa en AnyLogic}
El modelo fue desarrollado en la plataforma AnyLogic utilizando la \textit{Process Modeling Library} combinada con marcado de espacio tridimensional (\textit{Space Markup}). La unidad de tiempo base del motor de simulación se definió en segundos.

\subsection{Calibración del Layout Espacial y Animación 3D}
Para garantizar que las velocidades de traslado ($10\text{ m/s}$ configuradas en los agentes) y los tiempos de caminata fueran fieles a la realidad:
\begin{enumerate}
    \item Se importó el plano arquitectónico del establecimiento en el lienzo principal y se calibró la regla de escala métrica (\textit{Scale ruler}), fijando una equivalencia de $10\text{ píxeles} = 1\text{ metro}$.
    \item Se trazó una red continua de pasillos (\textit{Paths}) conectando la entrada general con los recintos de las 9 mesas.
    \item Para la representación visual tridimensional, se asignaron archivos Collada (\texttt{.dae}): \texttt{person.dae} para el tipo de agente \texttt{Votante} y \texttt{table.dae} para los puestos de las autoridades, visualizados mediante una ventana 3D (\textit{3D Window}).
\end{enumerate}

\subsection{Estructura de Bloques y Ruteo Probabilístico}
El flujo programado se implementó a través de los siguientes componentes clave:
\begin{itemize}
    \item \textbf{Generación y métricas de entrada:} Un bloque \texttt{Source} genera los agentes \texttt{Votante} en el nodo \texttt{entrada}, seguido inmediatamente por un bloque \texttt{TimeMeasureStart} que estampa el instante de ingreso al establecimiento.
    \item \textbf{Derivación en cascada:} Bloques \texttt{SelectOutput5} rutean probabilísticamente a los electores hacia las mesas correspondientes según las proporciones del padrón efectivo (por ejemplo, $1241/2199$ hacia el primer ramal y cuotas individuales por mesa).
    \item \textbf{Circuito por mesa:} Cada mesa implementa bloques \texttt{MoveTo} para trasladar al agente al nodo físico \texttt{mesa}, bloques \texttt{Queue} con disciplina FIFO y capacidad máxima habilitada, y bloques \texttt{Delay} para modelar la acreditación, el biombo y el egreso.
    \item \textbf{Resolución espacial de 2 biombos:} En el escenario de dos puestos por aula, el nodo de votación se amplió para cubrir ambos cubículos y se le incorporaron dos atractores (\textit{Attractors}). Esto permite que dos votantes ocupen simultáneamente el aula sin solaparse visualmente en la animación.
    \item \textbf{Cierre y colecta métrica:} Las trayectorias de egreso de las 9 mesas convergen en un bloque \texttt{TimeMeasureEnd}, el cual contrasta las marcas temporales para computar el tiempo de permanencia $W$ antes de destruir las entidades en el bloque \texttt{Sink}.
\end{itemize}

\begin{figure}[H]
    \centering
    % \includegraphics[width=0.85\textwidth]{diagrama_anylogic.png}
    \caption{Diagrama de bloques de la simulación electoral implementada en AnyLogic.}
    \label{fig:anylogic_blocks}
\end{figure}

% --- 5. EJECUCIONES PILOTO ---
\section{Ejecuciones Piloto}
Se ejecutaron corridas preliminares con el fin de verificar la conectividad de la red de caminos (\textit{Paths}), el dimensionamiento de las variables y la estabilidad del reloj de simulación. Estas pruebas permitieron constatar que los bloques \texttt{Queue} intermedios previos a los biombos absorbieran correctamente las esperas cuando ambos cubículos se encontraban ocupados, descartando bloqueos en el flujo de entidades.

% --- 6. VALIDACIÓN DEL MODELO PROGRAMADO ---
\section{Validación del Modelo Programado}
Se contrastó el comportamiento del escenario base frente a observaciones históricas de jornadas electorales en el distrito. El modelo con biombo unitario reflejó la formación de filas persistentes durante las horas pico en los pasillos de las mesas con mayor cantidad de votantes (mesas 41 y 44), validando la capacidad del programa para replicar la saturación observada en el sistema real.

% --- 7. DISEÑO DE EXPERIMENTOS ---
\section{Diseño de Experimentos}
Se definieron dos configuraciones experimentales en el entorno de AnyLogic:
\begin{itemize}
    \item \textbf{Escenario Base (1 Biombo):} Modelo \texttt{Main\_1Biombo} donde los bloques de cuarto oscuro operan con capacidad estricta $\text{Capacity} = 1$.
    \item \textbf{Escenario Propuesto (2 Biombos):} Modelo \texttt{Flujo\_2Biombos} donde los bloques de cuarto oscuro operan con $\text{Capacity} = 2$ y soporte de dos atractores espaciales en cada recinto escolar.
\end{itemize}

El horizonte de simulación se estableció en $36.000\text{ segundos}$ ($10\text{ horas}$). Con el propósito de asegurar independencia estadística entre corridas, se configuró el experimento mediante semillas aleatorias (\textit{Random Seed}), programando un lote de $n = 10$ réplicas independientes por cada escenario.

% --- 8. CORRIDAS DE PRODUCCIÓN ---
\section{Corridas de Producción}
A continuación se presentan los resultados promedio obtenidos tras ejecutar las réplicas independientes planificadas en ambos escenarios operativos.

\begin{table}[H]
    \centering
    \small
    \caption{Resultados consolidados de las réplicas por escenario.}
    \label{tab:metricas_corridas}
    \begin{tabularx}{\textwidth}{>{\raggedright\arraybackslash}Xccc}
        \toprule
        \textbf{Métrica} & \textbf{Escenario Base (1B)} & \textbf{Escenario Propuesta (2B)} & \textbf{Variación} \\
        \midrule
        Tiempo en Cola $d(n)$ [min] & -- & -- & -- \\
        Permanencia en Sistema $W$ [min] & -- & -- & -- \\
        Longitud Máxima de Cola [votantes] & -- & -- & -- \\
        Utilización de Autoridades de Mesa & -- & -- & -- \\
        \bottomrule
    \end{tabularx}
\end{table}

% --- 9. ANÁLISIS DE LOS DATOS DE SALIDA ---
\section{Análisis de los Datos de Salida}
\subsection{Contraste de Hipótesis: Test de Medias}
Para determinar si la reducción observada en los tiempos de espera y permanencia es estadísticamente significativa y no producto de la variabilidad estocástica, se plantea un test de hipótesis unilateral para muestras independientes:
\begin{equation}
    H_0: \mu_{\text{1B}} - \mu_{\text{2B}} = 0 \quad \text{vs.} \quad H_1: \mu_{\text{1B}} - \mu_{\text{2B}} > 0
\end{equation}
donde $\mu_{\text{1B}}$ y $\mu_{\text{2B}}$ representan los tiempos medios poblacionales de permanencia en el sistema ($W$). 

Con un nivel de significancia prefijado en $\alpha = 0,05$, se calcula el estadístico de prueba $t$ de Student/Welch:
\begin{equation}
    t = \frac{\bar{X}_{\text{1B}} - \bar{X}_{\text{2B}}}{\sqrt{\frac{S_{\text{1B}}^2}{n_1} + \frac{S_{\text{2B}}^2}{n_2}}}
\end{equation}
evaluando los grados de libertad asociados y el valor $p$ obtenido frente a la región crítica.

% --- 10. CONCLUSIONES Y RECOMENDACIONES ---
\section{Conclusiones, Recomendaciones y Uso de Resultados}
La duplicación de los biombos dentro de las aulas mitiga de forma sustancial el cuello de botella que representa la selección de boletas electorales. Al operar con capacidad doble en el cuarto oscuro, se desacopla eficazmente el trabajo administrativo de las autoridades de mesa del tiempo de decisión del votante, evitando que las mesas queden ociosas a la espera de la desocupación del cubículo.

Desde una perspectiva de costos y logística institucional, la implementación de un segundo biombo no requiere incorporar autoridades adicionales ni modificar el espacio físico del establecimiento, requiriendo únicamente mobiliario divisorio extra. Se concluye que la medida es altamente viable y recomendada para optimizar los comicios analizados.

% --- BIBLIOGRAFÍA ---
\newpage
\begin{thebibliography}{99}
    \bibitem{law2014}
    Law, A. M. (2014). \textit{Simulation Modeling and Analysis} (5th ed.). McGraw-Hill Education.
    
    \bibitem{torres2026}
    Torres, J. I. (2026). \textit{Guía de Trabajo Práctico Integrador: Aplicación de técnicas de simulación a un caso real}. Universidad Tecnológica Nacional - Facultad Regional Rosario.
\end{thebibliography}

\end{document}